\documentclass[12pt, a4paper]{article}

\usepackage[utf8]{inputenc}
\usepackage[T1]{fontenc}
\usepackage[spanish]{babel}
\usepackage{geometry}
\usepackage{hyperref}
\usepackage{listings}
\usepackage{xcolor}
\usepackage{enumitem}
\usepackage{booktabs}
\usepackage{titlesec}
\usepackage{fancyhdr}

\geometry{margin=2.5cm}

% Galio brand colors
\definecolor{galiogreen}{HTML}{28A745}
\definecolor{galiodark}{HTML}{343A40}
\definecolor{galiogray}{HTML}{ADB5BD}
\definecolor{codebg}{HTML}{F4F4F4}

% Code listing style
\lstdefinestyle{json}{
  backgroundcolor=\color{codebg},
  basicstyle=\ttfamily\footnotesize,
  breaklines=true,
  frame=single,
  rulecolor=\color{galiogray},
  captionpos=b,
  numbers=left,
  numberstyle=\tiny\color{galiogray},
}

\lstdefinestyle{bash}{
  backgroundcolor=\color{galiodark},
  basicstyle=\ttfamily\footnotesize\color{white},
  breaklines=true,
  frame=single,
  captionpos=b,
}

% Header / footer
\pagestyle{fancy}
\fancyhf{}
\lhead{\textcolor{galiogreen}{\textbf{Galio Electronics}}}
\rhead{US-PJ3-16}
\cfoot{\thepage}

% Section color
\titleformat{\section}{\large\bfseries\color{galiodark}}{}{0em}{}[\color{galiogreen}\titlerule]
\titleformat{\subsection}{\normalsize\bfseries\color{galiodark}}{}{0em}{}

\hypersetup{
  colorlinks=true,
  linkcolor=galiogreen,
  urlcolor=galiogreen,
}

% ─────────────────────────────────────────────
\begin{document}

\begin{titlepage}
  \centering
  \vspace*{3cm}
  {\color{galiogreen}\rule{\linewidth}{2pt}}\\[0.5cm]
  {\LARGE\bfseries US-PJ3-16}\\[0.3cm]
  {\Large Parametrizar endpoints y entornos\\(dev / stage / prod)}\\[0.3cm]
  {\color{galiogreen}\rule{\linewidth}{2pt}}\\[1cm]
  \begin{tabular}{rl}
    \textbf{Proyecto}   & Galio Mobile App (Flutter) \\
    \textbf{Equipo}     & MOBILE-FLUTTER, PLATFORM-TB, OPS-DEVOPS, QA-DOCS \\
    \textbf{Estado}     & Implementado \\
    \textbf{Fecha}      & Febrero 2026 \\
  \end{tabular}
  \vfill
  {\color{galiogray}\small Documento generado como parte del proceso de desarrollo interno.}
\end{titlepage}

% ─────────────────────────────────────────────
\section{Contexto y objetivo}

La aplicación mobile de Galio se comunica con una instancia privada de
ThingsBoard (TB) en \texttt{tb.galio.dev}. Antes de esta historia, la URL del
servidor estaba hardcodeada como \texttt{defaultValue} en
\texttt{app\_constants.dart}, lo que impedía cambiar de entorno sin modificar
código fuente.

\textbf{Objetivo:} eliminar todo hardcode de endpoints y proveer una
configuración limpia, versionable y segura para los entornos
\textit{dev}, \textit{stage} y \textit{prod}.

% ─────────────────────────────────────────────
\section{Estrategia elegida: \texttt{--dart-define-from-file}}

Flutter 3.7+ soporta pasar un archivo JSON en tiempo de compilación mediante
la flag \texttt{--dart-define-from-file}. Cada clave del JSON se convierte en
una variable de entorno accesible con \texttt{String.fromEnvironment()}, que
es exactamente el mecanismo que ya usaba \texttt{app\_constants.dart}.

\subsection{Por qué esta estrategia}

\begin{itemize}
  \item Builds directamente sobre el mecanismo existente --- sin cambios en la
        lógica de lectura de constantes.
  \item Un solo archivo JSON por entorno concentra toda la configuración.
  \item CI/CD solo cambia qué archivo JSON se pasa al build.
  \item Sin dependencias externas ni paquetes adicionales.
  \item Extensible: si en el futuro se necesitan Flutter Flavors (apps
        separadas por entorno en el mismo dispositivo), se agregan encima sin
        romper nada.
\end{itemize}

% ─────────────────────────────────────────────
\section{Archivos creados}

\subsection{Estructura de carpetas}

\begin{lstlisting}[style=bash]
config/
  config.example.json   <-- va al repositorio (plantilla sin secretos)
  dev.json              <-- .gitignore (emulador / local)
  stage.json            <-- .gitignore (staging)
  prod.json             <-- .gitignore (tb.galio.dev)
\end{lstlisting}

\subsection{Variables de configuración (\texttt{config.example.json})}

\begin{lstlisting}[style=json, caption=config/config.example.json]
{
  "thingsboardApiEndpoint": "https://your-tb-server.example.com",
  "thingsboardOAuth2CallbackUrlScheme": "dev.galio.app.auth",
  "appLinksUrlHost": "your-tb-server.example.com",
  "androidApplicationId": "dev.galio.app",
  "androidApplicationName": "Galio",
  "thingsboardIosAppSecret": "",
  "thingsboardAndroidAppSecret": "",
  "navigationType": "",
  "API_CALLS": "false",
  "VERBOSE": "false",
  "showAppVersion": "false"
}
\end{lstlisting}

\subsection{Diferencias entre entornos}

\begin{center}
\begin{tabular}{llccc}
  \toprule
  \textbf{Variable} & \textbf{dev} & \textbf{stage} & \textbf{prod} \\
  \midrule
  \texttt{thingsboardApiEndpoint} & \texttt{http://10.0.2.2:8080} & \texttt{https://tb-stage.galio.dev} & \texttt{https://tb.galio.dev} \\
  \texttt{androidApplicationName} & Galio Dev & Galio Stage & Galio \\
  \texttt{VERBOSE} & true & false & false \\
  \texttt{API\_CALLS} & true & false & false \\
  \texttt{showAppVersion} & true & true & false \\
  \bottomrule
\end{tabular}
\end{center}

\textbf{Nota:} \texttt{10.0.2.2} es la IP del host desde el emulador Android.
Para iOS Simulator usar \texttt{localhost}.

% ─────────────────────────────────────────────
\section{Cambios en código}

\subsection{Eliminación del hardcode (\texttt{app\_constants.dart})}

Se eliminó el \texttt{defaultValue} de \texttt{thingsBoardApiEndpoint}:

\begin{lstlisting}[style=json, language={}, caption=lib/constants/app\_constants.dart (antes)]
static const thingsBoardApiEndpoint = String.fromEnvironment(
  'thingsboardApiEndpoint',
  defaultValue: 'https://tb.galio.dev',  // <- hardcode eliminado
);
\end{lstlisting}

Si no se pasa configuración, \texttt{fromEnvironment} devuelve cadena vacía,
lo que activa la validación de arranque descrita a continuación.

\subsection{Validación de arranque (\texttt{main.dart})}

Se agregó una guarda al inicio de \texttt{main()} antes de cualquier
inicialización:

\begin{lstlisting}[style=json, language={}, caption=lib/main.dart]
if (ThingsboardAppConstants.thingsBoardApiEndpoint.isEmpty) {
  FlutterNativeSplash.remove();
  runApp(const _MissingConfigApp());
  return;
}
\end{lstlisting}

\texttt{\_MissingConfigApp} muestra una pantalla de error legible con el
comando exacto que el desarrollador debe ejecutar para arrancar correctamente.
No lanza excepciones ni produce crasheos silenciosos.

% ─────────────────────────────────────────────
\section{Seguridad: \texttt{.gitignore}}

Los archivos \texttt{dev.json}, \texttt{stage.json} y \texttt{prod.json} se
excluyeron del repositorio:

\begin{lstlisting}[style=bash]
# .gitignore
config/dev.json
config/stage.json
config/prod.json
\end{lstlisting}

Solo \texttt{config.example.json} se versiona. Nunca deben committearse
secretos ni URLs privadas.

% ─────────────────────────────────────────────
\section{Comandos de uso diario}

\begin{lstlisting}[style=bash]
# Copiar plantilla la primera vez
cp config/config.example.json config/dev.json
# (editar config/dev.json con los valores reales)

# Desarrollo (emulador Android)
flutter run --dart-define-from-file=config/dev.json

# Staging
flutter run --dart-define-from-file=config/stage.json

# Produccion
flutter run --dart-define-from-file=config/prod.json

# Build release
flutter build apk --dart-define-from-file=config/prod.json
flutter build ipa --dart-define-from-file=config/prod.json
\end{lstlisting}

% ─────────────────────────────────────────────
\section{CI/CD: inyección de config desde secretos}

Los archivos de config no están en el repo, por lo que el pipeline debe
generarlos antes del build a partir de secretos del repositorio:

\begin{lstlisting}[style=bash, caption=Ejemplo GitHub Actions]
- name: Crear config de produccion
  run: echo '${{ secrets.PROD_CONFIG_JSON }}' > config/prod.json

- name: Build APK
  run: flutter build apk --dart-define-from-file=config/prod.json
\end{lstlisting}

\texttt{PROD\_CONFIG\_JSON} es un secret del repositorio con el contenido
completo del JSON de producción.

% ─────────────────────────────────────────────
\section{Checklist de conectividad}

Antes de reportar un error de conexión, verificar este checklist:

\begin{enumerate}[label=\textcolor{galiogreen}{\checkmark}]
  \item \textbf{URL accesible} --- El dispositivo/emulador puede alcanzar el
        host configurado en \texttt{thingsboardApiEndpoint}.
  \item \textbf{Puerto abierto} --- Puerto \texttt{443} (HTTPS) o \texttt{8080}
        (HTTP) no bloqueado por firewall.
  \item \textbf{TLS válido} --- Certificado no expirado y firmado por CA
        reconocida (sin autofirmados, salvo config extra).
  \item \textbf{Health endpoint} --- \texttt{GET /api/noauth/health} devuelve
        \texttt{\{"status":"UP"\}}.
  \item \textbf{OAuth scheme} --- Si se usa SSO, el
        \texttt{thingsboardOAuth2CallbackUrlScheme} coincide con lo registrado
        en la instancia TB.
  \item \textbf{App links host} --- \texttt{appLinksUrlHost} coincide con el
        dominio real del servidor.
  \item \textbf{Archivo JSON cargado} --- La app arranca sin mostrar la
        pantalla de ``Configuración no encontrada''.
\end{enumerate}

% ─────────────────────────────────────────────
\section{Criterios de aceptación verificados}

\begin{center}
\begin{tabular}{lc}
  \toprule
  \textbf{Criterio} & \textbf{Estado} \\
  \midrule
  App conecta a instancia privada con config limpia & \textcolor{galiogreen}{\checkmark} \\
  No hay hardcodes críticos en código               & \textcolor{galiogreen}{\checkmark} \\
  Checklist de conectividad existe y sirve          & \textcolor{galiogreen}{\checkmark} \\
  \bottomrule
\end{tabular}
\end{center}

\end{document}
